% problem_id: S046-thm:qinv-notprop
% source_id: S046 (HoTT Book)
% source_file: equivalences.tex
% statement_locator: equivalences.tex lines 120--122
% proof_locator: equivalences.tex lines 123--137
% env: thm
% pairing: tight (verified)
% status: complete (statement + author proof verbatim + reason + locators)
%
\begin{thm}\label{thm:qinv-notprop}
  There exist types $A$ and $B$ and a function $f:A\to B$ such that $\qinv(f)$ is not a mere proposition.
\end{thm}

\begin{proof}
  It suffices to exhibit a type $X$ such that $\prd{x:X} (x=x)$ is not a mere proposition.
  Define $X\defeq \sm{A:\type} \brck{\bool=A}$, as in the proof of \cref{thm:no-higher-ac}.
  It will suffice to exhibit an $f:\prd{x:X} (x=x)$ which is unequal to $\lam{x} \refl{x}$.

  Let $a \defeq (\bool,\bproj{\refl{\bool}}) : X$, and let $q:a=a$ be the path corresponding to the nonidentity equivalence $e:\eqv\bool\bool$ defined by $e(\bfalse)\defeq\btrue$ and $e(\btrue)\defeq\bfalse$.
  We would like to apply \cref{lem:autohtpy} to build an $f$.
  By definition of $X$, equalities in subset types (\cref{subsec:prop-subsets}), and univalence, we have $\eqv{(a=a)}{(\eqv{\bool}{\bool})}$, which is a set, so~\ref{item:autohtpy1} holds.
  Similarly, by definition of $X$ and equalities in subset types we have~\ref{item:autohtpy2}.
  Finally, \cref{ex:eqvboolbool} implies that every equivalence $\eqv\bool\bool$ is equal to either $\idfunc[\bool]$ or $e$, so we can show~\ref{item:autohtpy3} by a four-way case analysis.

  Thus, we have $f:\prd{x:X} (x=x)$ such that $f(a) = q$.
  Since $e$ is not equal to $\idfunc[\bool]$, $q$ is not equal to $\refl{a}$, and thus $f$ is not equal to $\lam{x} \refl{x}$.
  Therefore, $\prd{x:X} (x=x)$ is not a mere proposition.
\end{proof}
